\documentclass[11pt,a4paper]{article}
\usepackage[margin=1in]{geometry}
\usepackage{hyperref}
\usepackage{enumitem}
\usepackage{graphicx}
\usepackage{xcolor}
\usepackage{booktabs}
\usepackage{tabularx}
\usepackage{array}

\makeatletter
\newcommand{\BvrSetLabInstructor}[1]{\gdef\Bvr@LabInstructor{#1}}
\newcommand{\Bvr@LabInstructor}{Mr. Jeelani Asif}
\newcommand{\BvrLabInstructorValue}{\Bvr@LabInstructor}
\makeatother

\usepackage{titling}
\pretitle{\begin{center}\bfseries\fontsize{18bp}{18bp}\selectfont}
\makeatletter
\newcommand{\BvrSubtitle}[1]{%
  \posttitle{%
    \par\end{center}
    \begin{center}\large#1\end{center}
    \vskip 0.5em}%
}
\makeatother

\title{DeadlineDesk: A Campus Web System for Placement\\
Rounds and Academic Assignment Deadlines}

\BvrSetLabInstructor{Mr. Jeelani Asif}

\BvrSubtitle{%
  Project Proposal for UCS503P \\
  Submitted to: \BvrLabInstructorValue \\ \bfseries
  Thapar Institute of Engineering and Technology%
}

\author{
  Aryan Gupta, CSED%
  \thanks{Roll No: \texttt{1024030764}; Email: \texttt{agupta15\_be24\@thapar.edu}}
  \and
  Aksh Goyal, CSED%
  \thanks{Roll No: \texttt{1024030766}; Email: \texttt{agoyal2\_be24\@thapar.edu}}
  \and
  Naveen Bansal, CSED%
  \thanks{Roll No: \texttt{1024030767}; Email: \texttt{nbansal3\_be24\@thapar.edu}}
}

\date{18 August 2026}

\begin{document}
\maketitle

\begin{abstract}
\noindent
Placement application windows and assignment cutoffs at TIET are often
communicated through WhatsApp forwards, shared drives, and scattered LMS
notices. Students miss deadlines; teaching assistants apply late rules
inconsistently. This proposal presents \textbf{DeadlineDesk}, a role-based
web system that combines a placement track (companies, rounds, document
checklists, reminders) with an academic dropbox (submission, late policy,
similarity stub, TA grading). Success is measured by on-time completion and
correct late-policy behaviour, with a working prototype targeted for Week~7.
\end{abstract}

\section{Introduction / Problem Statement}

Campus academic and placement calendars place concurrent pressure on students.
Placement cells publish company-specific application windows and document
requirements; departments publish assignment due times with late policies.
In practice, these obligations travel through informal channels: forwarded
messages, spreadsheets, and drive folders. Context is lost, reminders are
uneven, and last-minute rushes are common.

The gap is the lack of a single, enforceable system for two related deadline
workflows. On the placement side, students rarely have a structured checklist
and reminder trail for each round. On the academic side, submissions may bypass
the LMS; late acceptance is decided manually; grading status is hard to track.
Nearby course projects do not address this: CaseRoom focuses on interview
practice and scoring; College Mentorship focuses on long-term mentor matching;
generic campus-resource tools do not model round checklists or late-policy
submission states.

DeadlineDesk closes the gap with a deployable web product: shared
authentication and roles; a Placement Track for companies, rounds, document
checklists, and reminders; and an Academic Dropbox for assignments, file
upload, automated late evaluation, a lightweight similarity check stub, and
TA grading. The problem is locally relevant, actionable with mature web
frameworks, and suitable for measurable evaluation within one semester,
consistent with UCS503P expectations for time-to-value and continuous
integration.

\section{Objectives}

\textbf{Primary goal.}
Build a campus web system that makes placement-round readiness and assignment
submission outcomes visible, enforceable, and testable by Week~7, then refine
design artefacts and quality for the end-semester evaluation.

\textbf{Specific objectives.}
\begin{enumerate}[leftmargin=0.7cm,itemsep=2pt]
  \item Implement role-based authentication for Student, TA/Faculty, and
        Placement Admin, verified by login tests for each role.
  \item By Week~7, deliver one end-to-end placement path: create
        company/round $\rightarrow$ document checklist $\rightarrow$
        reminder log entry ($T{-}24$h).
  \item By Week~7, deliver one end-to-end academic path: publish assignment
        with late policy $\rightarrow$ submit $\rightarrow$ automatic late
        flag $\rightarrow$ similarity stub score $\rightarrow$ TA grade.
  \item Maintain 100\% agreement between automated late flags and the stated
        policy on a fixed black-box suite of at least five automated tests.
  \item Keep the similarity check as an explicit stub (not a research NLP
        system); run CI (build and tests) and publish project documentation
        via the course MkDocs Pages pipeline.
\end{enumerate}

\section{Methodology}

This is an engineering methodology: architecture, module development,
integration, and validation.

\subsection{System architecture and technology stack}
We shall use a layered web architecture:
\begin{itemize}[leftmargin=0.7cm,itemsep=1pt]
  \item \textbf{Presentation:} responsive web UI for student, TA, and admin
        portals.
  \item \textbf{Application:} API for authentication, CRUD operations, status
        transitions, and reminder scheduling hooks.
  \item \textbf{Domain:} shared entities (User, Role, Deadline, Document,
        Reminder, StatusHistory) and module entities (Company, Round,
        ChecklistItem; Assignment, Submission, Grade).
  \item \textbf{Persistence:} relational database (PostgreSQL; SQLite
        acceptable in the early prototype) with indexed deadline and status
        queries.
\end{itemize}
Planned stack (final choice by team delivery speed): Django or Next.js,
GitHub Actions for CI, MkDocs with GitHub Pages from the UCS503P template,
and Pytest or Jest for tests. The design relies on mature components rather
than novel algorithms.

\subsection{Module development}
\begin{enumerate}[leftmargin=0.7cm,itemsep=2pt]
  \item \textbf{Authentication and roles.} Session or JWT authentication with
        role claims and route guards.
  \item \textbf{Placement Track.} An administrator publishes a company and
        round with open/close times. Round states:
        draft $\rightarrow$ published $\rightarrow$ open $\rightarrow$
        closed $\rightarrow$ archived.
        Students complete required checklist items; readiness requires all
        required items to be complete.
  \item \textbf{Academic Dropbox.} A TA sets the due time and late policy
        (reject / accept-with-penalty / grace-window). Late evaluation is a
        function of submission time, due time, and policy. Submission states:
        draft $\rightarrow$ submitted $\rightarrow$ under-review $\rightarrow$
        graded (with an optional late flag).
  \item \textbf{Reminders and audit.} Persist scheduled and sent reminder
        records; outbound email is a later increment.
\end{enumerate}

\subsection{Integration, CI/CD, and validation}
Work proceeds in weekly increments: build, test, and demonstrate thin features.
CI runs linting and unit tests on each push; documentation deploys through the
template Pages pipeline; a staging URL supports the Week~7 demonstration.

\textbf{Evaluation.}
Primary metric: on-time completion rate (OTCR)---the fraction of tracked
deadlines where the required student action occurs before cutoff, taken from
database timestamps. Secondary metrics: reminder effectiveness; full agreement
of the late-policy test suite; median TA turnaround time. Pilot plan: 8--15
student users; staff roles may be simulated by the team if needed.

\textbf{Risks and mitigations.}
Two-module scope is controlled by a shared core and thin Week~7 paths.
Reminder delivery on free hosting starts as a database log. The similarity
component is labelled as a stub in the SRS. Official college email addresses
are recorded on the team declaration sheet.

\textbf{Scalability.}
Pagination, indexes, and clear module boundaries support growth in users and
deadlines. Deployment targets common web hosting with a managed database and
CI/CD.

\section{Timeline}

\begin{center}
\small
\begin{tabular}{@{}p{1.6cm}p{6.2cm}p{5.0cm}@{}}
\toprule
\textbf{Week} & \textbf{Task} & \textbf{Deliverable} \\
\midrule
W3--W4 & Problem definition, proposal, presentation, repository setup &
  Proposal PDF, slides, Pages site \\
W5--W6 & SRS draft, use-case and class sketches, backlog, auth scaffold &
  SRS v0.1, CI green \\
\textbf{W7} & Both thin end-to-end paths and automated tests &
  \textbf{Prototype (MST)} \\
W8 & Improvement plan from prototype feedback &
  Improvement note \\
W9--W10 & Mid-semester buffer; defect fixes &
  Stable staging build \\
W11--W12 & Sequence and state diagrams; additional tests; UI polish &
  Design pack v1 \\
W13 & Buffer (holiday sandwich) &
  Catch-up \\
W14--W15 & Second prototype polish &
  Prototype 2 \\
W16 & Improvements over second prototype &
  Updated demonstration \\
W17 & Final prototype, seminar, final report &
  \textbf{EST delivery} \\
\bottomrule
\end{tabular}
\end{center}

\vspace{0.3em}
\noindent\textit{Buffer:} Weeks~13 and the mid-semester period absorb
slippage. If needed, email reminders and dashboards are deferred before core
status machines.

\section{Resources / Budget}

\textbf{Human resources.}
\begin{itemize}[leftmargin=0.7cm,itemsep=1pt]
  \item \textbf{Aryan Gupta (1024030764):} architecture, CI/CD, placement
        module.
  \item \textbf{Aksh Goyal (1024030766):} academic dropbox module and late-policy
        tests.
  \item \textbf{Naveen Bansal (1024030767):} user interface flows,
        documentation site, demonstration script.
\end{itemize}
Role emphasis may adjust after Week~7 according to measured velocity.

\textbf{Material resources.} Personal laptops; GitHub (repository, Actions,
Pages); free-tier hosting and database for staging; \LaTeX\ for reports. No
specialised laboratory hardware is required.

\textbf{Budget.} No external funding is required. Software is free and
open-source or available through student GitHub access. Optional printing of a
one-page handout is incidental.

\section{Summary}

DeadlineDesk is a focused software engineering project: two deadline
workflows, clear roles, explicit status machines, continuous integration, and
measurable on-time and policy-correctness metrics. It matches UCS503P
selection heuristics---local relevance, available engines, deployability, and
evaluation---and is attainable for a three-member team within the semester.

\end{document}
